\subsubsection{Global Convergence}
\label{app:pduot-global}

We give formal versions of the two conclusions in Theorem~\ref{prop:pduot-convergence}.
Throughout this section, fix the active support $\mathcal N$ and abbreviate $\mathbf c=\mathbf c_{\mathcal N}$.
Convergence refers to the iteration continued without finite stopping, with exact proximal evaluations.

\begin{theorem}[Global convergence of \subsolver]
\label{thm:pduot-global}
Assume that $\mathcal N$ covers every source and target node, the reference masses and penalty coefficients are strictly positive, and the transport costs are finite.
Fix positive diagonal base steps satisfying $\chi:=\|\overline{\mathbf D}_y^{1/2}\mathbf A_{\mathcal N}\overline{\mathbf D}_x^{1/2}\|_2<1$.
In Algorithm~\ref{alg:pduot}, use $\mathbf D_x=w_j\overline{\mathbf D}_x$ and $\mathbf D_y=\overline{\mathbf D}_y/w_j$ in epoch $j$, with $0<w_{\min}\leq w_j\leq w_{\max}<\infty$ and the summable log-weight budget~\eqref{eq:pduot-weight-budget}.
Restarts may occur after any map evaluation and set the next anchor to that mapped point.
Then, from any finite initial anchor, the iterates and mapped points converge to a common restricted KKT point $\mathbf z^\infty=(\mathbf x^\infty,\mathbf y^\infty)$.
The primal projection multipliers also converge, to the restricted dual slack $\mathbf c+\mathbf A_{\mathcal N}^{\top}\mathbf y^\infty$.
\end{theorem}
Within a restart epoch, Algorithm~\ref{alg:pduot} applies reflected Halpern iteration:
\begin{equation}
\label{eq:halpern}
\mathbf z^{k+1}=\frac{1}{k+2}\mathbf z^{\mathrm a}+\frac{k+1}{k+2}(2\pduotmap\mathbf z^k-\mathbf z^k).
\end{equation}
Here $\mathbf z^{\mathrm a}$ is the fixed epoch anchor and $\mathbf z^0=\mathbf z^{\mathrm a}$.
For a positive definite matrix $\mathbf{P}$, write $\|\mathbf{z}\|_\mathbf{P}=(\mathbf{z}^\top \mathbf{P}\mathbf{z})^{1/2}$ and $\operatorname{dist}_\mathbf{P}(\mathbf{z},S)=\inf_{u\in S}\|\mathbf{z}-\mathbf{u}\|_\mathbf{P}$.
Define the restricted KKT mapping
\begin{equation}
\label{eq:pduot-kkt-map}
\mathcal M(\mathbf{x},\mathbf{y})=\begin{pmatrix}\mathbf{c}+\mathbf{A}_{\mathcal N}^\top \mathbf{y}+N_{\mathbb R_+^{|\mathcal N|}}(\mathbf{x})\\ \nabla G^*(\mathbf{y})-\mathbf{A}_{\mathcal N}\mathbf{x}\end{pmatrix},
\end{equation}
where $N_{\mathbb R_+^{|\mathcal N|}}$ is the convex normal-cone mapping.
The zeros of $\mathcal M$ are precisely the restricted KKT points, with $\mathbf{y}=-(\mathbf{f},\mathbf{g})$.
The mapping $\mathcal M$ is maximally monotone: $\partial F\times\partial G^*$ is maximally monotone, and the remaining linear map $(\mathbf{x},\mathbf{y})\mapsto(\mathbf{A}_{\mathcal N}^\top \mathbf{y},-\mathbf{A}_{\mathcal N}\mathbf{x})$ is skew-symmetric and defined everywhere.

For fixed positive diagonal steps, define the PDHG metric
\begin{equation*}
\mathbf{P}=\begin{pmatrix}\mathbf{D}_x^{-1}&-\mathbf{A}_{\mathcal N}^\top\\-\mathbf{A}_{\mathcal N}&\mathbf{D}_y^{-1}\end{pmatrix}.
\end{equation*}

\begin{proof}
\textit{KKT existence and the fixed-metric map.}
The restricted primal objective is continuous on the nonnegative orthant and coercive: both transported marginal totals equal $\|\mathbf x\|_1$, so the KL penalties grow superlinearly in $\|\mathbf x\|_1$, whereas the finite transport costs contribute a linear term.
It therefore attains a minimum $\mathbf x^\star$.
Suppose its source marginal $r_i^\star$ is zero.
Choose an incident active edge $e=(i,j)$ and add flow $t>0$ on that edge.
The source KL change divided by $t$ is $\rho_a(\log(t/a_i)-1)$, which tends to $-\infty$ as $t\downarrow0$.
The cost quotient is finite, while the target KL quotient tends either to the finite value $\rho_b\log(q_j^\star/b_j)$ if $q_j^\star>0$, or to $-\infty$ otherwise.
The total quotient is therefore negative for sufficiently small $t$, contradicting optimality.
The target argument is symmetric, so $\mathbf{r}^\star,\mathbf{q}^\star>0$.
Setting $\mathbf y^\star=\nabla G(\mathbf A_{\mathcal N}\mathbf x^\star)$, the primal first-order optimality condition and conjugacy give $0\in\mathcal M(\mathbf x^\star,\mathbf y^\star)$, so the restricted KKT solution set is nonempty.
The step-size condition makes $\mathbf P$ positive definite by the Schur complement.
For $\widehat{\mathbf{z}}=T\mathbf{z}$, the two proximal optimality conditions in~\eqref{eq:pduot-map} are
\begin{equation*}
\mathbf{D}_x^{-1}(\mathbf{x}-\widehat{\mathbf{x}})-\mathbf{A}_{\mathcal N}^\top \mathbf{y}\in\partial F(\widehat{\mathbf{x}}),\qquad
\mathbf{D}_y^{-1}(\mathbf{y}-\widehat{\mathbf{y}})+\mathbf{A}_{\mathcal N}(2\widehat{\mathbf{x}}-\mathbf{x})=\nabla G^*(\widehat{\mathbf{y}}).
\end{equation*}
Rearranging gives
\begin{equation}
\label{eq:pduot-resolvent}
\mathbf{P}(\mathbf{z}-T\mathbf{z})\in\mathcal M(T\mathbf{z}),\qquad T=(\mathbf{P}+\mathcal M)^{-1}\mathbf{P}.
\end{equation}
Monotonicity applied to the images of two points $\mathbf{z},\mathbf{u}$ yields
\begin{equation*}
\|T\mathbf{z}-T\mathbf{u}\|_\mathbf{P}^2\leq\langle T\mathbf{z}-T\mathbf{u},\mathbf{z}-\mathbf{u}\rangle_\mathbf{P}.
\end{equation*}
Thus $T$ is firmly nonexpansive, $R:=2T-\mathbf{I}$ is nonexpansive, and $\operatorname{Fix}R=\operatorname{Fix}T=\mathcal M^{-1}(0)$.
Equation~\eqref{eq:halpern}, initialized at $\mathbf{z}^0=\mathbf{z}^{\mathrm a}$, is the Halpern iteration for $R$ with coefficients $1/(k+2)$.
The Halpern convergence theorem~\citep{wittmann1992approximation} and the residual estimate of \citet[Theorem~2.1]{lieder2021convergence}, applied to $R$ in the $\mathbf P$ inner product, give convergence to a fixed point and
\begin{equation*}
\|\mathbf{z}^k-R\mathbf{z}^k\|_\mathbf{P}\leq\frac{2\operatorname{dist}_\mathbf{P}(\mathbf{z}^{\mathrm a},\operatorname{Fix}R)}{k+1}.
\end{equation*}
Since $\mathbf{z}^k-R\mathbf{z}^k=2(\mathbf{z}^k-T\mathbf{z}^k)$, we obtain
\begin{equation}
\label{eq:pduot-halpern-bound}
\|\mathbf z^k-T\mathbf z^k\|_{\mathbf P}\leq\frac{\operatorname{dist}_{\mathbf P}(\mathbf z^{\mathrm a},\operatorname{Fix}T)}{k+1}.
\end{equation}
Continuity of $T$ gives convergence of the mapped points to the same limit in a final infinite epoch.

The degree-based steps in~\eqref{eq:pduot-metric} satisfy the required norm condition without estimating a spectral norm.
Indeed, Cauchy--Schwarz at each node gives, for any edge vector $\mathbf{h}$,
\begin{equation*}
\|\overline{\mathbf{D}}_y^{1/2}\mathbf{A}_{\mathcal N}\overline{\mathbf{D}}_x^{1/2}\mathbf{h}\|_2^2
=\frac{\zeta^2}{2}\sum_v\frac{1}{d_v}\left(\sum_{e\ni v}\mathbf{h}_e\right)^2
\leq\frac{\zeta^2}{2}\sum_v\sum_{e\ni v}\mathbf{h}_e^2
=\zeta^2\|\mathbf{h}\|_2^2.
\end{equation*}
Each edge has two endpoints, and the reciprocal factors $w$ in $\mathbf{D}_x$ and $\mathbf{D}_y$ cancel in the scaled incidence matrix.
Hence its norm is at most $\zeta<1$ for every fixed $w>0$.
\par\smallskip\noindent\textit{Metric changes under the weight budget.}
Write $s=\log w$ and
\begin{equation*}
\mathbf P(s)=\begin{pmatrix}e^{-s}\overline{\mathbf D}_x^{-1}&-\mathbf A_{\mathcal N}^{\top}\\-\mathbf A_{\mathcal N}&e^s\overline{\mathbf D}_y^{-1}\end{pmatrix},\qquad T_s=(\mathbf P(s)+\mathcal M)^{-1}\mathbf P(s).
\end{equation*}
Set $c_0=(1-\chi)^{-1}$.
If $\mathbf D(s)$ is the block-diagonal part of $\mathbf P(s)$, then $\mathbf P(s)\succeq(1-\chi)\mathbf D(s)$ and, for every nonzero $\mathbf h$,
\begin{equation*}
\left|\frac{d}{ds}\log(\mathbf h^\top\mathbf P(s)\mathbf h)\right|\leq\frac{\mathbf h^\top\mathbf D(s)\mathbf h}{\mathbf h^\top\mathbf P(s)\mathbf h}\leq c_0.
\end{equation*}
Integration gives
\begin{equation}
\label{eq:pduot-metric-comparison}
e^{-c_0|t-s|}\mathbf P(s)\preceq\mathbf P(t)\preceq e^{c_0|t-s|}\mathbf P(s).
\end{equation}
For an integer $J\geq1$, the budget in~\eqref{eq:pduot-weight-budget} satisfies
\begin{equation*}
\sum_{j=1}^{\infty}|s_j-s_{j-1}|\leq\sum_{j=1}^{\infty}b_j\leq LJ\left(1+\frac{1}{p-1}\right)<\infty.
\end{equation*}
Consequently $s_j$ converges, and $\mathbf P_j:=\mathbf P(s_j)$ converges to a positive definite matrix.
The compact weight interval also gives constants $0<\underline\lambda\leq\overline\lambda<\infty$ with $\underline\lambda\mathbf I\preceq\mathbf P_j\preceq\overline\lambda\mathbf I$ for all $j$.

\par\smallskip\noindent\textit{Convergence under mapped-point restarts.}
Suppose first that infinitely many restarts occur.
Index epochs by $j\geq0$, with anchor $\mathbf a_j$, fixed map $T_j:=T_{s_j}$, and final Halpern iterate $\mathbf t_j$.
The next anchor is $\mathbf a_{j+1}=T_j\mathbf t_j$.
Both $T_j$ and $2T_j-\mathbf I$ fix every $\mathbf z^\star\in\mathcal S:=\mathcal M^{-1}(0)$ and are nonexpansive in $\mathbf P_j$.
Induction in~\eqref{eq:halpern} therefore bounds every iterate and mapped point in epoch $j$ by the anchor's $\mathbf P_j$-distance to $\mathbf z^\star$.
Firm nonexpansiveness at the final iterate gives, with $e_j=\|\mathbf t_j-T_j\mathbf t_j\|_{\mathbf P_j}$,
\begin{equation}
\label{eq:pduot-epoch-descent}
\|\mathbf a_{j+1}-\mathbf z^\star\|_{\mathbf P_j}^2+e_j^2\leq\|\mathbf a_j-\mathbf z^\star\|_{\mathbf P_j}^2.
\end{equation}
Set $r_j=e^{c_0b_{j+1}}$, $\Gamma_j=\prod_{i=0}^{j-1}r_i$, and $E_j=\|\mathbf a_j-\mathbf z^\star\|_{\mathbf P_j}^2$.
Equations~\eqref{eq:pduot-metric-comparison}--\eqref{eq:pduot-epoch-descent} imply
\begin{equation*}
\frac{E_{j+1}}{\Gamma_{j+1}}\leq\frac{E_j}{\Gamma_j}-\frac{e_j^2}{\Gamma_j},\qquad 1\leq\Gamma_j\leq\Gamma_\infty<\infty.
\end{equation*}
Thus the anchors are bounded, $E_j$ has a limit for every solution $\mathbf z^\star$, and $\sum_je_j^2<\infty$.
The resolvent inclusion gives $\mathbf P_j(\mathbf t_j-\mathbf a_{j+1})\in\mathcal M(\mathbf a_{j+1})$.
Its left-hand side tends to zero by $e_j\to0$ and the uniform spectral bounds.
Closedness of the maximal monotone graph puts every cluster point of the anchors in $\mathcal S$.
Choose one such point $\mathbf z^\infty$.
The energy $E_j$ relative to this point has a limit and a subsequence tending to zero, so the whole anchor sequence converges to $\mathbf z^\infty$.
The within-epoch distance bound and uniform spectral bounds give chronological convergence of all Halpern iterates and mapped points to the same limit.
If only finitely many restarts occur, the last epoch has fixed weight, and the Halpern convergence theorem gives the same conclusion.

The primal projection multiplier satisfies
\begin{equation*}
\boldsymbol\lambda=\mathbf c+\mathbf A_{\mathcal N}^{\top}\mathbf y-\mathbf D_x^{-1}(\mathbf x-\widehat{\mathbf x}).
\end{equation*}
Since $\mathbf D_x^{-1}$ is uniformly bounded, these multipliers also converge, to $\mathbf c+\mathbf A_{\mathcal N}^{\top}\mathbf y^\infty$.
\end{proof}

\subsubsection{Local Linear Convergence}
\label{app:pduot-restart}

\begin{theorem}[Local linear convergence of \subsolver]
\label{thm:pduot-local}
Under the assumptions of Theorem~\ref{thm:pduot-global}, let $\mathcal R_j(\mathbf z)=\|\mathbf z-T_j\mathbf z\|_{\mathbf P_j}$ be the fixed-point residual in epoch $j$.
Assume that restart checks have a uniformly bounded gap of at most $h$ PDHG evaluations, including the first check of each epoch.
For a fixed $\vartheta\in(0,1)$, restart whenever $\mathcal R_j(\mathbf z)\leq\vartheta\mathcal R_j(\mathbf a_j)$, evaluating the anchor baseline under the current metric; other restart triggers are allowed.
Then the iterates, mapped points, and primal projection multipliers converge locally R-linearly in the number of PDHG evaluations.
Specifically, there exist $C>0$, $\gamma\in(0,1)$, and $N_0$ such that, at evaluation $N\geq N_0$,
\begin{equation*}
\|\mathbf z_N-\mathbf z^\infty\|_2+\|\widehat{\mathbf z}_N-\mathbf z^\infty\|_2+\|\boldsymbol\lambda_N-\boldsymbol\lambda^\infty\|_2\leq C\gamma^{N-N_0}.
\end{equation*}
\end{theorem}

\begin{proof}

\textit{Local KKT error bound.}
\label{app:pduot-sharpness}
Introduce the transported marginal vector $\mathbf m=(\mathbf r,\mathbf q)$ and write the restricted problem as
\begin{equation*}
\min_{\mathbf m,\mathbf x}\;G(\mathbf m)+F(\mathbf x)\quad\text{s.t.}\quad\mathbf A_{\mathcal N}\mathbf x-\mathbf m=0.
\end{equation*}
This problem satisfies the structured polyhedricity assumption of \citet[Assumptions~1.2--1.3]{yuan2020discerning}.
Indeed, the extended-valued KL function $G$ is proper, closed, and convex, is differentiable on the positive orthant, and has
\begin{equation*}
\nabla^2G(\mathbf m)=\operatorname{diag}(\rho_a/\mathbf r,\rho_b/\mathbf q)\succ0.
\end{equation*}
Thus $G$ is strongly convex on every compact convex subset of its subdifferential domain, its gradient is locally Lipschitz on the positive orthant, and the gradient norm diverges when approaching a finite boundary point.
In the notation of that assumption, take the smooth representation matrix to be the identity and the additional linear term to be zero.
The function $F$ is convex piecewise linear-quadratic, since it is a linear function plus the indicator of the nonnegative orthant.
Consequently, \citet[Theorem~59]{yuan2020discerning} gives metric subregularity of the lifted KKT mapping at every lifted KKT point, without rank or strict-complementarity assumptions.

To transfer this bound to~\eqref{eq:pduot-kkt-map}, define, for $\mathbf m>0$,
\begin{equation*}
\mathcal K(\mathbf m,\mathbf x,\mathbf y)=\begin{pmatrix}\nabla G(\mathbf m)-\mathbf y\\ \mathbf c+\mathbf A_{\mathcal N}^{\top}\mathbf y+N_{\mathbb R_+^{|\mathcal N|}}(\mathbf x)\\ \mathbf m-\mathbf A_{\mathcal N}\mathbf x\end{pmatrix}.
\end{equation*}
Up to signs and ordering of the residual blocks, this is the KKT mapping in the cited theorem.
Write $\mathcal S=\mathcal M^{-1}(0)$ and $\widetilde{\mathcal S}=\mathcal K^{-1}(0)$, and fix $\mathbf z^\star=(\mathbf x^\star,\mathbf y^\star)\in\mathcal S$.
The projection of $\widetilde{\mathcal S}$ onto $(\mathbf x,\mathbf y)$ is exactly $\mathcal S$.
For $\mathbf z=(\mathbf x,\mathbf y)$ near $\mathbf z^\star$, set $\mathbf m(\mathbf y)=\nabla G^*(\mathbf y)$.
This positive vector depends continuously on $\mathbf y$, and $\nabla G(\mathbf m(\mathbf y))=\mathbf y$.
The lifted error bound therefore yields, for some $\kappa>0$ and a neighborhood $\mathbf U$ of $\mathbf z^\star$,
\begin{equation}
\label{eq:pduot-kkt-error-bound}
\begin{aligned}
\operatorname{dist}_2(\mathbf z,\mathcal S)&\leq\operatorname{dist}_2((\mathbf m(\mathbf y),\mathbf x,\mathbf y),\widetilde{\mathcal S})\\
&\leq\kappa\operatorname{dist}_2(0,\mathcal K(\mathbf m(\mathbf y),\mathbf x,\mathbf y))\\
&=\kappa\operatorname{dist}_2(0,\mathcal M(\mathbf z)),\qquad \mathbf z\in\mathbf U.
\end{aligned}
\end{equation}
This bound measures distance to the entire solution set and allows nonunique optimal couplings, cycles in the positive-flow support, and edges with both zero flow and zero slack.

\par\smallskip\noindent\textit{Uniform fixed-point sharpness.}
Global convergence places all sufficiently late iterates in the KKT error-bound neighborhood of $\mathbf z^\infty$ from~\eqref{eq:pduot-kkt-error-bound}.
The uniform spectral bounds transfer that error bound to every late-epoch map, giving a common constant $\alpha_0>0$ such that
\begin{equation}
\label{eq:pduot-uniform-sharpness}
\alpha_0\operatorname{dist}_{\mathbf P_j}(\mathbf z,\mathcal S)\leq\|\mathbf z-T_j\mathbf z\|_{\mathbf P_j}
\end{equation}
near the limit.
Indeed, the resolvent inclusion and~\eqref{eq:pduot-kkt-error-bound} imply
\begin{equation*}
\operatorname{dist}_{\mathbf P_j}(T_j\mathbf z,\mathcal S)\leq\sqrt{\overline\lambda}\kappa\|\mathbf P_j(\mathbf z-T_j\mathbf z)\|_2\leq\kappa\overline\lambda\|\mathbf z-T_j\mathbf z\|_{\mathbf P_j}.
\end{equation*}
The triangle inequality gives $\alpha_0=(1+\kappa\overline\lambda)^{-1}$.
Nonexpansiveness and uniform norm equivalence ensure that all $T_j\mathbf z$ remain in the error-bound neighborhood for $\mathbf z$ sufficiently near $\mathbf z^\infty$.

\par\smallskip\noindent\textit{Contraction across epochs.}
Write $d_j=\operatorname{dist}_{\mathbf P_j}(\mathbf a_j,\mathcal S)$.
Nonexpansiveness gives $\mathcal R_j(T_j\mathbf t_j)\leq e_j$.
Applying~\eqref{eq:pduot-uniform-sharpness} at $\mathbf a_{j+1}=T_j\mathbf t_j$, and taking a nearest solution to $\mathbf a_j$ in~\eqref{eq:pduot-epoch-descent}, yields
\begin{equation*}
\operatorname{dist}_{\mathbf P_j}(\mathbf a_{j+1},\mathcal S)^2+e_j^2\leq d_j^2,\qquad \alpha_0\operatorname{dist}_{\mathbf P_j}(\mathbf a_{j+1},\mathcal S)\leq e_j.
\end{equation*}
With $q_0=(1+\alpha_0^2)^{-1/2}<1$, metric comparison then gives
\begin{equation*}
d_{j+1}\leq e^{c_0b_{j+1}/2}q_0d_j.
\end{equation*}
Since $b_j\to0$, for any $q\in(q_0,1)$ the right-hand factor is at most $q$ in all sufficiently late epochs.
Moreover $\|\mathbf a_{j+1}-\mathbf a_j\|_2\leq2d_j/\sqrt{\underline\lambda}$ by the within-epoch distance bound.
Summing these geometrically decreasing increments proves R-linear convergence to the actual limit, uniformly for all iterates and mapped points within each late epoch, when indexed by epoch.

\par\smallskip\noindent\textit{Rate per PDHG evaluation.}
The Halpern estimate~\eqref{eq:pduot-halpern-bound} and uniform sharpness give
\begin{equation*}
\mathcal R_j(\mathbf z^{j,k})\leq\frac{d_j}{k+1}\leq\frac{\mathcal R_j(\mathbf a_j)}{\alpha_0(k+1)}.
\end{equation*}
Thus the practical sufficient-reduction condition
\begin{equation*}
\mathcal R_j(\mathbf z^{j,k})\leq\vartheta\mathcal R_j(\mathbf a_j)
\end{equation*}
holds once $k+1\geq1/(\alpha_0\vartheta)$.
Every sufficiently late epoch therefore has at most $H=\lceil1/(\alpha_0\vartheta)\rceil+h$ PDHG evaluations; additional restart triggers can only shorten it.
A final infinite epoch from a nonsolution is impossible, because its residual tends to zero and eventually passes the sufficient-reduction test.
If an anchor is already a solution, every subsequent iterate equals it.
Otherwise, after $N$ evaluations in the tail at least $\lfloor N/H\rfloor$ epochs have completed, giving a bound $Cq^{\lfloor N/H\rfloor}$ for the iterates and mapped points.
The multiplier formula in the global proof and bounded $\mathbf D_x^{-1}$ give the same rate for $\boldsymbol\lambda_N$.
\end{proof}

The degree-based preconditioner in Appendix~\ref{app:pduot-controls} satisfies the global step-size condition, and its restart rule uses $\vartheta=0.2$ with periodic checks.
The moderate-reduction and epoch-length triggers are additional restarts permitted by both theorems.
